\documentclass[11pt]{article}
\usepackage[T1]{fontenc}
\usepackage[margin=1in]{geometry}
\usepackage{amsmath,amssymb,amsthm,mathtools}
\usepackage{graphicx}
\usepackage{float}
\usepackage{microtype}
\usepackage[hidelinks]{hyperref}

\newtheorem{theorem}{Theorem}
\newtheorem{lemma}[theorem]{Lemma}
\newtheorem{corollary}[theorem]{Corollary}
\newtheorem{remark}[theorem]{Remark}
\newcommand{\R}{\mathbb R}
\newcommand{\M}{\mathcal M}
\newcommand{\Mt}{\widetilde{\mathcal M}}
\newcommand{\N}{\mathcal N}
\newcommand{\Nt}{\widetilde{\mathcal N}}
\newcommand{\vphi}{\varphi}
\newcommand{\tvphi}{\widetilde{\varphi}}
\newcommand{\tr}{\operatorname{tr}}
\newcommand{\supp}{\operatorname{supp}}

\title{Power Tail Decay Gives the Noncommutative Weak Endpoint\\
       in the Correct Kernel Orientation}
\author{Partial resolution for arXiv:1702.06536}
\date{13 August 2026}

\begin{document}
\maketitle

\begin{center}
\fbox{\parbox{0.91\linewidth}{\textbf{Status.}
Candidate partial resolution, likely valid; human review required.  The
quantitative tail condition gives weak type $(1,1)$ for the operator when it
is imposed in the input variable, including the source's
$\M'\cap\Mt$-valued kernels.  The one-sided formula printed in the source is
in the opposite orientation and gives the endpoint for the Hilbert adjoint.
Under the natural two-sided reading, both endpoints follow.}}
\end{center}

\section{The source question and the orientation issue}

Let $(\Mt,\widetilde\tau)$ be a semifinite von Neumann algebra and let
$\M\subset\Mt$ be a von Neumann subalgebra on which $\widetilde\tau$ remains
semifinite.  Put
\[
 \N=L_\infty(\R^d)\,\bar\otimes\,\M,
 \qquad
 \Nt=L_\infty(\R^d)\,\bar\otimes\,\Mt.
\]
Their natural tensor traces are denoted
$\vphi=\int_{\R^d}dx\otimes\tau$ and
$\tvphi=\int_{\R^d}dx\otimes\widetilde\tau$.
The source studies an $L_2(\N)\to L_2(\Nt)$ operator
\begin{equation}\label{eq:T}
 Tf(x)=\int_{\R^d} k(x,y)f(y)\,dy,
 \qquad k(x,y)\in \M'\cap\Mt,
\end{equation}
with
\begin{equation}\label{eq:size}
 \|k(x,y)\|_{\Mt}\le C_0|x-y|^{-d}.
\end{equation}
Its Remark~1.4 asks whether a quantitative integral tail can replace
pointwise Lipschitz estimates:

\begin{figure}[H]
\centering
\includegraphics[page=4,width=0.93\linewidth,trim=72 318 72 38,clip]
  {source_paper.pdf}
\caption{The exact question in Remark~1.4 of \cite{source}.  Notice that the
displayed difference varies the first argument of $k$.}
\label{fig:source}
\end{figure}

For \eqref{eq:T}, cancellation of a bad input on a cube $Q$ produces
\[
 k(x,y)-k(x,c_Q),
\]
so weak type $(1,1)$ for $T$ naturally asks for regularity in the
\emph{second}, or input, variable.  The displayed condition in
Figure~\ref{fig:source} varies the \emph{first}, or output, variable.  The
distinction cannot be erased by renaming dummy variables.

We prove a stronger power-tail theorem in the input orientation, extend the
later scalar-kernel endpoint argument to the source's commutant-valued
setting, and then state exactly what the printed condition implies.

\section{Statement of the result}

For $\beta>0$, say that $k$ satisfies the input-tail condition
$\mathrm{Tail}_{\rm in}(A,\beta)$ when, for all $y,z\in\R^d$ and $\alpha>2$,
\begin{equation}\label{eq:tailin}
 \int_{|x-y|>\alpha|y-z|}
 \|k(x,y)-k(x,z)\|_{\Mt}\,dx
 \le A\alpha^{-\beta}.
\end{equation}
The exponent need not equal the dimension.

\begin{theorem}[Input-tail endpoint]\label{thm:main}
Assume \eqref{eq:size}, $\mathrm{Tail}_{\rm in}(A,\beta)$ for some
$\beta>0$, and boundedness $T:L_2(\N)\to L_2(\Nt)$.  Then $T$ is of weak
type $(1,1)$: there is a finite constant
\[
 C=C\bigl(d,\beta,C_0,A,\|T\|_{2\to2}\bigr)
\]
such that for every $f\in L_1(\N)$ and $\lambda>0$,
\begin{equation}\label{eq:weak}
 \lambda\,\tvphi\!\left(
   \mathbf 1_{(\lambda,\infty)}(|Tf|)\right)
 \le C\|f\|_{L_1(\N)}.
\end{equation}
\end{theorem}

In particular, the source's decay $\alpha^{-d}$ is more than enough in the
input orientation.  A useful slightly more general formulation replaces
$\alpha^{-\beta}$ in \eqref{eq:tailin} by a decreasing modulus $\omega(\alpha)$
such that
\begin{equation}\label{eq:sqrtdini}
 \sum_{j\ge0}\sqrt{\omega(c_d2^j)}<\infty.
\end{equation}

\section{Proof intuition}

The later paper \cite{ccap} found a Calder\'on--Zygmund decomposition that
avoids pseudolocalization.  Its endpoint theorem asks for a square-integral
of the input-variable kernel difference on every remote annulus.  The source's
tail estimate gives an $L_1$ bound there.  The size estimate gives an
$L_\infty$ bound there.  Multiplying the two bounds gives the required
$L_2$ estimate, and taking a square root explains why the Dini condition is
on $\sqrt\omega$.

The published proof is written for scalar kernels and notes a central-valued
extension.  For the source's more general kernels, every kernel value commutes
with the input algebra.  This is exactly what the row/column factorization in
the proof needs: $|K|^2$ can be replaced by $\|K\|^2\mathbf1$ under the trace,
with no change to the Cuculescu decomposition.

\section{Tail decay implies annular \texorpdfstring{$L_2$}{L2}--H\"ormander regularity}

Fix an odd dilation factor $D_d$ large enough that, for a cube $Q$ of side
$\ell(Q)$, $y,z\in Q$, and $x\notin D_dQ$, the distances $|x-y|$ and
$|x-z|$ are comparable to $|x-c_Q|$.  Let
\[
 E_j(Q)=\left\{x\notin D_dQ:
  2^jD_d\ell(Q)\le |x-c_Q|<2^{j+1}D_d\ell(Q)\right\},
 \qquad j\ge0.
\]

\begin{lemma}[Tail-to-square upgrade]\label{lem:tail2}
If \eqref{eq:size} and \eqref{eq:tailin} hold, then
\begin{equation}\label{eq:H2}
 \sup_Q\sum_{j\ge0}\sup_{y\in Q}
 \left(
  |E_j(Q)|\int_{E_j(Q)}
  \|k(x,y)-k(x,c_Q)\|_{\Mt}^{,2}\,dx
 \right)^{1/2}<\infty.
\end{equation}
More precisely, its left side is at most
$C_{d,\beta}\sqrt{AC_0}$.  The same conclusion holds with a modulus
$\omega$ satisfying \eqref{eq:sqrtdini}.
\end{lemma}

\begin{proof}
Fix $Q,j,y$, and write $z=c_Q$.  On $E_j(Q)$,
\[
 |x-y|\simeq_d 2^j\ell(Q),\qquad
 |x-z|\simeq_d 2^j\ell(Q),\qquad |y-z|\lesssim_d\ell(Q).
\]
After increasing $D_d$ once, the annulus lies in the tail region of
\eqref{eq:tailin} with $\alpha=c_d2^j>2$.  Hence
\begin{equation}\label{eq:L1shell}
 \int_{E_j(Q)}\|k(x,y)-k(x,z)\|_{\Mt}\,dx
 \le C_dA2^{-j\beta}.
\end{equation}
The size condition gives
\begin{equation}\label{eq:Linfshell}
 \sup_{x\in E_j(Q)}\|k(x,y)-k(x,z)\|_{\Mt}
 \le C_dC_0,2^{-jd}\ell(Q)^{-d}.
\end{equation}
Since $|E_j(Q)|\le C_d2^{jd}\ell(Q)^d$, multiplying
\eqref{eq:L1shell} and \eqref{eq:Linfshell} yields
\[
 |E_j(Q)|\int_{E_j(Q)}
   \|k(x,y)-k(x,z)\|_{\Mt}^2\,dx
 \le C_dAC_0,2^{-j\beta}.
\]
Its square root is summable because $\beta>0$.  Replacing
$2^{-j\beta}$ by $\omega(c_d2^j)$ proves the last assertion.
\end{proof}

The fixed dilation $D_d$ only changes dimensional constants in the standard
exceptional projection.  Equivalently, one may retain the dilation used in
\cite{ccap}; size controls the finitely many extra shells, and
Lemma~\ref{lem:tail2} controls the remote ones.

\section{The commutant-valued endpoint argument}

\begin{lemma}[Extension of the later endpoint theorem]\label{lem:commutant}
Let $T$ have the form \eqref{eq:T}, be bounded
$L_2(\N)\to L_2(\Nt)$, and have a kernel in $\M'\cap\Mt$ satisfying
\eqref{eq:H2}.  Then \eqref{eq:weak} holds.
\end{lemma}

\begin{proof}
We audit the proof of Theorem~A in \cite{ccap}.  At height $\lambda$, its new
noncommutative Calder\'on--Zygmund decomposition writes
\[
 f=g+b_{\rm d}+b_{\rm off}
\]
using Cuculescu projections in $\N$, and constructs a dilated projection
$\zeta\in\N$ with
\[
 \vphi(\mathbf1-\zeta)
 \lesssim_d\lambda^{-1}\|f\|_1.
\]
The good term uses only Chebyshev, the $L_2$ bound for $T$, and
$\|g\|_2^2\lesssim_d\lambda\|f\|_1$, so it is unchanged when the range is
$\Nt$.  Terms carrying $\mathbf1-\zeta$ are controlled by the same support
projection estimate and the fact that $\widetilde\tau|_{\M}=\tau$.

For a bad diagonal piece supported on $Q$, cancellation gives
\[
 \zeta(x)T b_{{\rm d},Q}(x)\zeta(x)
 =\zeta(x)\int_Q
 [k(x,y)-k(x,c_Q)]b_{{\rm d},Q}(y)\,dy\,\zeta(x)
\]
outside the dilated cube.  The trace-ideal inequality
$\|ab\|_1\le\|a\|_\infty\|b\|_1$ and the $L_1$ consequence of
\eqref{eq:H2} give the published diagonal estimate verbatim.

Only the off-diagonal piece uses the scalar nature of the kernel in the
printed proof.  Write
\[
 K_Q(x,y)=k(x,y)-k(x,c_Q).
\]
The factors $\pi_Qf(y)\pi_Q$ and all Cuculescu projections lie in $\M$, while
$K_Q(x,y)\in\M'\cap\Mt$.  They therefore commute, and positivity gives the
exact replacement for the scalar square used in \cite{ccap}:
\begin{equation}\label{eq:keytrace}
 \widetilde\tau\!\left(
 |K_Q(x,y)|^2\pi_Qf(y)\pi_Q\right)
 \le \|K_Q(x,y)\|_{\Mt}^{,2}
     \tau\!\left(\pi_Qf(y)\pi_Q\right).
\end{equation}
The row/column Cauchy--Schwarz factorization, shell integration, and
\eqref{eq:H2} now give the same bound as in the scalar proof:
\[
 \sum_Q\|\zeta T b_{{\rm off},Q}\zeta\|_1
 \lesssim_d \mathfrak H_2(k)\|f\|_1,
\]
where $\mathfrak H_2(k)$ denotes the left side of \eqref{eq:H2}.  Summing the
good, exceptional, diagonal, and off-diagonal distribution estimates proves
\eqref{eq:weak}.  No step asks that $K_Q$ be central; commutation with $\M$
is sufficient.
\end{proof}

Theorem~\ref{thm:main} follows immediately from
Lemmas~\ref{lem:tail2} and \ref{lem:commutant}.

\section{What the formula printed in the source proves}

Define the output-tail condition $\mathrm{Tail}_{\rm out}(A,\beta)$ by
\begin{equation}\label{eq:tailout}
 \int_{|x-y|>\alpha|x-x'|}
 \|k(x,y)-k(x',y)\|_{\Mt}\,dy
 \le A\alpha^{-\beta}
 \qquad(\alpha>2).
\end{equation}
This is exactly the orientation printed in Figure~\ref{fig:source} when
$\beta=d$.

\begin{corollary}[Adjoint and two-sided conclusions]\label{cor:orientation}
Assume the setup of Theorem~\ref{thm:main}.
\begin{enumerate}
\item If \eqref{eq:tailout} holds, the Hilbert-space adjoint
$T^\dagger:L_2(\Nt)\to L_2(\N)$, restricted to $L_1(\N)$, is weak type
$(1,1)$.
\item If both \eqref{eq:tailin} and \eqref{eq:tailout} hold, then $T$ and
$T^\dagger|_{L_1(\N)}$ are weak type $(1,1)$.
\item When $\M=\Mt$, $T^\dagger$ is the usual integral adjoint with kernel
$k^\dagger(u,v)=k(v,u)^*$.
\end{enumerate}
\end{corollary}

\begin{proof}
Because the ambient trace is tracial, there is a trace-preserving conditional
expectation $E_\M:\Mt\to\M$.  On separated supports the restriction of the
Hilbert adjoint has kernel
\[
 h(u,v)=E_\M\bigl(k(v,u)^*\bigr).
\]
Conditional expectation is contractive in operator norm, and bimodularity
shows that $h(u,v)\in\M\cap\M'=\mathcal Z(\M)$.  Condition
\eqref{eq:tailout} for $k$ is condition \eqref{eq:tailin} for $h$, with the
same constants.  Also $\|T^\dagger\|_{2\to2}=\|T\|_{2\to2}$.  Apply
Theorem~\ref{thm:main} to the adjoint.  The second assertion combines this
with Theorem~\ref{thm:main} for $T$, and the third is immediate.
\end{proof}

\section{Scope, verification, and literature relation}

The result answers the weak-endpoint question affirmatively under the natural
two-sided replacement of the source's preceding two Lipschitz estimates.  It
also shows that \emph{any} positive power tail, not only $\alpha^{-d}$, is
sufficient in the input orientation.  It does not prove that the original
pseudolocalization estimate survives: the later decomposition was designed
to avoid that estimate.

The literal one-sided formula \eqref{eq:tailout} only supplies the endpoint
for the adjoint.  Weak $(1,1)$ for $T^\dagger$ plus the $L_2$ bound gives
strong bounds for $T$ above $2$ by interpolation and duality, but it does not
reach weak $L_1$ for $T$.  A proof-grade counterexample to the remaining
one-sided implication was not found; no such implication is claimed here.

Hong--Lai--Xu \cite{hlx} prove an endpoint under a summable annular $L_2$
condition for scalar kernels tensored with the identity.  Cadilhac,
Conde-Alonso, and Parcet \cite{ccap} give the decomposition used above and an
annular $L_2$--H\"ormander theorem, noting extensions to arbitrary
noncommutative input spaces and central-valued kernels.  The new deductions
here are the square-Dini tail criterion, the explicit orientation correction,
and the commutant-valued range extension matching \cite{source}.  Novelty is
provisional; eight distinct proof and counterexample upgrades are recorded in
the accompanying attempt log.

\begin{thebibliography}{9}
\bibitem{source}
L. Cadilhac,
\emph{Weak boundedness of Calder\'on--Zygmund operators on noncommutative
$L_1$-spaces}, J. Funct. Anal. 274 (2018), 769--796;
arXiv:1702.06536.

\bibitem{hlx}
G. Hong, X. Lai, and B. Xu,
\emph{Maximal singular integral operators acting on noncommutative
$L_p$-spaces}, arXiv:2009.03827.

\bibitem{ccap}
L. Cadilhac, J. M. Conde-Alonso, and J. Parcet,
\emph{Spectral multipliers in group algebras and noncommutative
Calder\'on--Zygmund theory}, arXiv:2105.05036.
\end{thebibliography}

\end{document}
